\begin{center}
\begin{tabular}{@{}lll@{}}
\toprule
\textbf{\tAction} & \textbf{\tKeyboard} & \textbf{\tJoystick}\\
\midrule
Andar a izquierda / derecha & \key{$\leftarrow$} \key{$\rightarrow$} & izquierda / derecha\\
Entrar por una puerta o doblar una esquina & \key{$\uparrow$} & arriba\\
Lanzar un beso & \key{$\downarrow$} & abajo\\
Abrir la barra de iconos & \key{ESPACIO} & bot\'on (cualquiera de los dos)\\
\addlinespace
En la barra: elegir icono / confirmar & \key{$\leftarrow$} \key{$\rightarrow$} / \key{ESPACIO} & izquierda / derecha / bot\'on\\
En un men\'u: elegir l\'inea / confirmar & \key{$\uparrow$} \key{$\downarrow$} / \key{ESPACIO} & arriba / abajo / bot\'on\\
Cerrar un mensaje & \key{ESPACIO} & bot\'on\\
\bottomrule
\end{tabular}
\end{center}

\begin{itemize}
\item Para entrar en un edificio o doblar una esquina, ponte \textbf{de cara} a la puerta o a la esquina y pulsa arriba: Andy se gira de espaldas y entra.
\item El joystick funciona en el \textbf{puerto 1 y en el 2}.
\item Un mensaje se queda en pantalla hasta que pulsas \key{ESPACIO} o el bot\'on, para que te d\'e tiempo a leerlo. Mientras tanto el reloj sigue corriendo y la hora del panel se va actualizando.
\item No hay tecla de pausa.
\end{itemize}
